\paragraph{Removing the common entries in the sieve.}
There is a separate deduction behind the exponent \(2403n^9f\) in
the bound for \(Z\). Let \(A\) be the finite set of possible
nonzero entries, integral units outside the finite places used in
the sieve. Write \(A^i\) for the possible first \(i\) entries of
a head. Let \(J_i(p,q)\) be its level-\(i\) kernel, with all
unwritten entries padded by one and all entries taken with their
prescribed local scales. For an entry \(y\) of the local product,
put
\[
 F_i(y)=\sum_{p,q\in A^i} J_{i+1}((p,y),(q,y)),\qquad
 D_{i+1}=\sum_{t\in A} F_i(t),\quad D_0=\Re J_0((),()) .
\]
The diagonal kernel on a common entry is positive semidefinite,
so \(F_i\) is real and nonnegative. Assume the local integration
identity \(\int F_i=\Re\sum_{p,q}J_i(p,q)\) and let \(w_i(p)\ge0\)
be weights. More precisely, suppose the following estimates hold
in the measures giving Poisson constant \(P_i=C_{{\rm Pois},i}\):
\begin{enumerate}
\item One-sided Poisson summation applies to \(F_i\), with total
absolute value of its nonzero Fourier coefficients at most
\(\epsilon(\sum_{p\in A^i}w_i(p))^2\) in the dual sum.
\item For heads \(p,q\in A^i\) with distinct last entries,
\( |J_i(p,q)|\le\epsilon w_i(p)w_i(q)\).
\item \((\sum_{p\in A^i}w_i(p))^2\le
  \Gamma\prod_{j<i}P_j\) for \(0\le i<n\).
\end{enumerate}
The first estimate, with the zero frequency evaluated by the
integration identity, bounds \(D_{i+1}\) by
\(P_i(\Re\sum_{p,q}J_i(p,q)+\epsilon(\sum_p w_i(p))^2)\).
In this unrestricted head sum, the equal-last-entry contribution is
\(D_i\); the remaining pairs contribute at most
\(\epsilon(\sum_p w_i(p))^2\) by the second estimate.
At \(i=0\) there are no remaining pairs. Thus
\[
 \frac{D_{i+1}}{\prod_{j\le i}P_j}
 \le \frac{D_i}{\prod_{j<i}P_j}+2\epsilon\Gamma .
\]
The top sum restricted to allowable last entries satisfies \(Z\le D_n\)
by nonnegativity. Iteration gives
\(Z\le(\prod_{j<n}P_j)(D_0+2n\epsilon\Gamma)\).
The fitted kernel bounds needed here are
\[
 D_0\le \frac{(1+b^{-10})^f}{\delta}
           \frac{1+b^{-6}}{b^{6k|T|}},\qquad
 \epsilon=C_{\rm dec}\,b^{2002n^9f},\qquad
 \Gamma=b^{400n^9f}.
\]
Since \(2n\le b^{n^9 f}\), these bounds prove the stated sieve
inequality. They isolate the mass at level zero and the two
different uses of decay. For unequal last entries \(t,s\), the
finite frequency is \((t^{-1}-s^{-1})/M\), a nonzero element
integral outside the chosen places when \(1/M\) is integral there.
A pointwise decay bound at those frequencies supplies item 2.
Section~\ref{moment:sieve} proves the weight bound and sums an explicit
majorant over the full dual lattice. The actual kernels must satisfy
the local domination, integration, support and positivity hypotheses there.
